\begin{table}[htb!]
\centering
\caption{Evaluation of the modelled damage distribution against
FEMA-verified housing damage (primary) and the Huang et al.\ (2025)
building sample (secondary). For each storm, the scaled weighted units
affected (WUA) per county is rank-correlated (Spearman $r_s$, exact
two-sided $p$) with the FEMA-verified real property damage of owner
households from the OpenFEMA Individual Assistance records, over the
counties present in both sources ($n$); correlations are shown for
$n \geq 4$. Detection is the share of the storm's total FEMA-verified
damage, summed over all declared counties, that falls in counties with
modelled damage. The Huang columns repeat the comparison against the
number of assessed damaged buildings per county. The pooled row
correlates within-storm shares across all matched county-storm pairs.
Hermine produced no modelled damage (detection 0\%). Nate is not part
of the analysis (no reconstructed wind field).}
\label{tab:damage_distribution}
\begin{tabular}{lrrrrrrr}
\hline
 & \multicolumn{4}{c}{FEMA IA verified damage} & \multicolumn{3}{c}{Huang et al.\ building sample} \\
Storm & $n$ & $r_s$ & $p$ & Detection & $n$ & $r_s$ & $p$ \\
\hline
Katrina (2005) & 16 & $+0.34$ & 0.192 & 95\% & -- & -- & -- \\
Isaac (2012) & 1 & -- & -- & 16\% & 1 & -- & -- \\
Hermine (2016) & -- & -- & -- & 0\% & -- & -- & -- \\
Harvey (2017) & 9 & $+0.80$ & 0.010 & 4\% & 6 & $-0.03$ & 0.957 \\
Irma (2017) & 4 & $-0.60$ & 0.400 & 37\% & 4 & $-0.40$ & 0.600 \\
Michael (2018) & 29 & $+0.75$ & $<0.001$ & 99\% & 8 & $+0.93$ & $<0.001$ \\
Laura (2020) & 12 & $+0.65$ & 0.022 & 92\% & 7 & $+0.82$ & 0.025 \\
Ida (2021) & 18 & $+0.40$ & 0.097 & 96\% & -- & -- & -- \\
Ian (2022) & 8 & $+0.81$ & 0.015 & 82\% & -- & -- & -- \\
Pooled (within-storm shares) & 96 & $+0.56$ & $<0.001$ & -- & 25 & $+0.34$ & 0.101 \\
\hline
\end{tabular}
\end{table}
